\documentclass[11pt]{article}
\usepackage[a4paper,margin=1in]{geometry}
\usepackage{fontspec}
\usepackage[boldfont=false]{xeCJK}
\setCJKmainfont{FandolSong-Regular.otf}
\usepackage{amsmath}
\usepackage{amssymb}
\usepackage{array}
\usepackage{booktabs}
\usepackage{graphicx}
\usepackage{adjustbox}
\usepackage{enumitem}
\setlist[itemize]{topsep=2pt,partopsep=0pt,parsep=0pt,itemsep=2pt,leftmargin=1.4em}
\setlist[enumerate]{topsep=2pt,partopsep=0pt,parsep=0pt,itemsep=2pt,leftmargin=1.6em}
\usepackage{parskip}
\usepackage{fvextra}
\fvset{breaklines=true,breakanywhere=true,fontsize=\small}
\setlength{\parindent}{0pt}

\begin{document}

\begin{center}
  {\LARGE\bfseries Electrochemistry}\\[0.35em]
  {\large A Level Chemistry}
\end{center}
\vspace{0.6em}

\section*{Oxidation number}

The \textbf{oxidation number} 氧化数 (also called the oxidation state) shows how many \textbf{electrons} 电子 an atom has lost or gained compared with the free element. You work it out using simple rules:

\begin{itemize}
\item an uncombined element has an oxidation number of $0$.

\item a simple ion has an oxidation number equal to its charge (so $\text{Mg}^{2+}$ is $+2$).

\item Group 1 is always $+1$, Group 2 is always $+2$.

\item hydrogen is $+1$ (but $-1$ in metal hydrides).

\item oxygen is $-2$ (but $-1$ in peroxides).

\item fluorine is always $-1$.

\item the oxidation numbers in a neutral compound add up to $0$; in an ion they add up to the charge.

\end{itemize}

A Roman numeral shows the size of the oxidation number of an element, for example iron(II) means $+2$ and manganese(VII) in $\text{KMnO}_4$ means $+7$.

\section*{Redox in terms of electrons}

A \textbf{redox} 氧化还原 reaction is one where electrons move from one species to another.

\begin{itemize}
\item \textbf{oxidation} 氧化 is the \textbf{loss} of electrons. The oxidation number goes \textbf{up}.

\item \textbf{reduction} 还原 is the \textbf{gain} of electrons. The oxidation number goes \textbf{down}.

\end{itemize}

A useful memory aid is \textbf{OIL RIG}: Oxidation Is Loss, Reduction Is Gain.

Oxidation and reduction always happen together, because the electrons lost by one species are gained by another. This \textbf{electron transfer} 电子转移 is why we call it a redox reaction.

\section*{Using oxidation numbers to balance equations}

Changes in oxidation number help you balance a redox equation:

\begin{enumerate}
\item find the element whose oxidation number \textbf{rises} (it is oxidised) and the one whose number \textbf{falls} (it is reduced).

\item the total rise must equal the total fall, because every electron lost is gained somewhere.

\item choose the ratio of the two species so the rise and fall match, then balance the rest of the equation.

\end{enumerate}

\section*{Disproportionation}

\textbf{Disproportionation} 歧化 is a special redox reaction in which the \textbf{same} element is both oxidised and reduced at the same time. For example, when chlorine reacts with cold water, some chlorine atoms are reduced (to $\text{Cl}^-$) and others are oxidised (to $\text{ClO}^-$).

\section*{Oxidising and reducing agents}

\begin{itemize}
\item an \textbf{oxidising agent} 氧化剂 takes electrons away from another species. In doing so, it is itself reduced.

\item a \textbf{reducing agent} 还原剂 gives electrons to another species. In doing so, it is itself oxidised.

\end{itemize}

So in any redox reaction, the oxidising agent causes the oxidation of the other species, while the reducing agent causes the reduction.


\end{document}
